\chapter{Equivalence of Propositions}
\signature

\section{Overview}

Chapter~2 introduced logical equivalence through truth tables. This chapter
turns equivalence into a working calculus: classifying formulas by their
satisfiability profile, simplifying them with the standard laws, and putting
them into \emph{normal forms} suited to mechanical processing.

\section{Tautology, contradiction, contingency revisited}

\begin{definition}
Let $\varphi$ be a formula over $n$ atomic letters, with $2^n$ valuations.
$\varphi$ is a \emph{tautology} if all $2^n$ valuations satisfy it, a
\emph{contradiction} if none does, and a \emph{contingency} otherwise.
$\varphi$ is \emph{satisfiable} when at least one valuation satisfies it ---
i.e.\ exactly when it is not a contradiction.
\end{definition}

\begin{example}
$(p \to q) \leftrightarrow (\neg q \to \neg p)$ is a tautology (the
contrapositive law); $(p \oplus q) \wedge (p \leftrightarrow q)$ is a
contradiction; $p \to (q \wedge r)$ is a contingency.
\end{example}

\begin{notebox}
$\varphi$ is a tautology if and only if $\neg\varphi$ is a contradiction.
Tautology checking and satisfiability checking are therefore two faces of the
same computational problem.
\end{notebox}

\section{The main equivalences}

The equivalence laws of Chapter~2 --- identity, domination, idempotence,
double negation, commutativity, associativity, distributivity, De~Morgan,
absorption, negation --- together with the implication identities
\[
p \to q \equiv \neg p \vee q, \qquad
\neg(p \to q) \equiv p \wedge \neg q, \qquad
p \leftrightarrow q \equiv (p \to q) \wedge (q \to p),
\]
form a complete working kit. Two further conditional equivalences are worth
memorising:
\[
p \to (q \to r) \;\equiv\; (p \wedge q) \to r
\qquad\text{(exportation)},
\]
\[
(p \to q) \wedge (p \to r) \;\equiv\; p \to (q \wedge r),
\qquad
(p \to r) \wedge (q \to r) \;\equiv\; (p \vee q) \to r .
\]

\begin{example}[Simplifying a proposition]
Simplify $\varphi = (p \wedge q) \vee (p \wedge \neg q) \vee (\neg p \wedge q)$.
\begin{align*}
\varphi &\equiv \bigl(p \wedge (q \vee \neg q)\bigr) \vee (\neg p \wedge q)
  && \text{distributivity}\\
&\equiv p \vee (\neg p \wedge q) && \text{negation, identity}\\
&\equiv (p \vee \neg p) \wedge (p \vee q) && \text{distributivity}\\
&\equiv p \vee q && \text{negation, identity.}
\end{align*}
Three two-literal clauses collapse to a single disjunction.
\end{example}

Such simplification is exactly what automated tools (SAT preprocessors,
circuit minimisers, symbolic algebra systems) perform on a large scale;
the Quine--McCluskey algorithm and Karnaugh maps systematise it for circuit
design.

\section{Disjunctive normal form}

\begin{definition}
A \emph{minterm} over atoms $p_1,\dots,p_n$ is a conjunction containing each
$p_i$ exactly once, either plain or negated. A formula is in \emph{disjunctive
normal form} (DNF) if it is a disjunction of conjunctions of literals; it is in
\emph{full} (canonical) DNF if every disjunct is a minterm.
\end{definition}

\begin{theorem}[DNF from the truth table]
Every satisfiable formula $\varphi$ is equivalent to the disjunction of the
minterms of the rows where $\varphi$ is true. Hence every formula has a full
DNF, unique up to order.
\end{theorem}

\begin{example}[XOR: ``for me or against me'']
The exclusive or is true on rows $(\TT,\FF)$ and $(\FF,\TT)$, so its full DNF
is
\[
p \oplus q \;\equiv\; (p \wedge \neg q) \vee (\neg p \wedge q).
\]
\end{example}

\begin{example}[A counting function]
Define $m(p,q,r)$ to be true when \emph{at least two} of $p,q,r$ are true
(the ``majority'' function). The true rows are
$(\TT,\TT,\TT), (\TT,\TT,\FF), (\TT,\FF,\TT), (\FF,\TT,\TT)$, giving the full
DNF
\[
m \equiv (p\wedge q\wedge r) \vee (p\wedge q\wedge \neg r)
\vee (p\wedge \neg q\wedge r) \vee (\neg p\wedge q\wedge r),
\]
which simplifies (by absorption-style grouping) to
$(p \wedge q) \vee (p \wedge r) \vee (q \wedge r)$.
\end{example}

\section{Other normal forms}

\begin{definition}
A formula is in \emph{negation normal form} (NNF) if negation appears only on
atoms, and in \emph{conjunctive normal form} (CNF) if it is a conjunction of
\emph{clauses}, each clause a disjunction of literals. The full CNF uses
\emph{maxterms} --- clauses containing every atom once --- one for each
\emph{false} row of the truth table.
\end{definition}

CNF is the input format of SAT solvers and of the resolution proof system;
DNF makes satisfiability trivial to read off but may be exponentially large.

\subsection{CNF \texorpdfstring{$\leftrightarrow$}{<->} DNF conversion}

The two directions are symmetric: distribute $\vee$ over $\wedge$ to go from
CNF to DNF, and $\wedge$ over $\vee$ for the converse.

\begin{example}
\[
(p \vee q) \wedge (\neg p \vee r)
\;\equiv\;
(p \wedge \neg p) \vee (p \wedge r) \vee (q \wedge \neg p) \vee (q \wedge r)
\;\equiv\;
(p \wedge r) \vee (\neg p \wedge q) \vee (q \wedge r),
\]
a CNF of two clauses becoming a DNF of three conjunctions (the contradictory
term vanishes). The term $q \wedge r$ is in fact absorbed by no other and must
stay; in general the blow-up can be exponential.
\end{example}

\section{Satisfiability research}

Deciding satisfiability of a CNF formula is the historic NP-complete problem
(Cook--Levin, 1971). Despite the exponential worst case, modern
conflict-driven clause-learning (CDCL) solvers routinely handle industrial
instances with millions of variables, powering hardware verification, planning
and dependency resolution. The special case 2-SAT (clauses of width $2$) is
solvable in linear time, while 3-SAT is already NP-complete --- a remarkably
sharp frontier.

\section{Summary}

Formulas split into tautologies, contradictions and contingencies; the
equivalence laws constitute an algebra for simplifying them; every formula has
a CNF and a DNF, readable from the false and true rows of its truth table
respectively; and satisfiability, trivial for DNF, is the canonical hard
problem for CNF.

\section*{Exercises}
\addcontentsline{toc}{section}{Exercises}

\begin{exercise}{\easy}
Classify as tautology, contradiction or contingency:
(a) $(p \wedge q) \to p$;\quad
(b) $p \leftrightarrow \neg p$;\quad
(c) $(p \vee q) \wedge \neg p$.
\end{exercise}

\begin{exercise}{\easy}
Verify by truth table the exportation law
$p \to (q \to r) \equiv (p \wedge q) \to r$.
\end{exercise}

\begin{exercise}{\medium}
Simplify using the laws, citing each step:
$\neg\bigl(\neg p \wedge \neg q\bigr) \wedge (p \vee \neg q)$.
\end{exercise}

\begin{exercise}{\medium}
Simplify $(p \to q) \wedge (p \to \neg q)$ to a formula with a single
connective.
\end{exercise}

\begin{exercise}{\medium}
Write the full DNF of $p \to \neg q$ (two atoms), directly from its truth
table.
\end{exercise}

\begin{exercise}{\medium}
Give the full CNF of $p \oplus q$, using the false rows of its truth table.
\end{exercise}

\begin{exercise}{\hard}
Let $f(p,q,r)$ be true exactly when an \emph{odd} number of $p, q, r$ are
true. Write its full DNF and show $f \equiv p \oplus q \oplus r$.
\end{exercise}

\begin{exercise}{\hard}
Convert $(p \wedge q) \vee (r \wedge s)$ into CNF, and explain why a CNF of a
DNF with $k$ conjunctions of two literals each may need $2^k$ clauses.
\end{exercise}

\begin{exercise}{\hard}
Show that the majority function
$m \equiv (p\wedge q) \vee (p\wedge r) \vee (q\wedge r)$ is also equivalent to
the CNF $(p \vee q) \wedge (p \vee r) \wedge (q \vee r)$.
\end{exercise}

\section*{Solutions to the Exercises}
\addcontentsline{toc}{section}{Solutions to the Exercises}

\begin{solution}{3.1}
(a) Tautology --- if $p \wedge q$ holds, $p$ holds. (b) Contradiction --- the
two sides always differ. (c) Contingency --- true for $p=\FF, q=\TT$, false
for $p=\TT$.
\end{solution}

\begin{solution}{3.2}
Both sides are false exactly on the row $(p,q,r) = (\TT,\TT,\FF)$: the left
side because $q \to r$ then fails with $p$ true; the right side because
$p \wedge q$ is true and $r$ false. All other rows make both sides true,
hence equivalence.
\end{solution}

\begin{solution}{3.3}
\begin{align*}
\neg(\neg p \wedge \neg q) \wedge (p \vee \neg q)
&\equiv (p \vee q) \wedge (p \vee \neg q) && \text{De Morgan, double negation}\\
&\equiv p \vee (q \wedge \neg q) && \text{distributivity}\\
&\equiv p \vee \FF \equiv p && \text{negation, identity.}
\end{align*}
\end{solution}

\begin{solution}{3.4}
$(p \to q) \wedge (p \to \neg q) \equiv p \to (q \wedge \neg q)
\equiv p \to \FF \equiv \neg p$.
\end{solution}

\begin{solution}{3.5}
$p \to \neg q$ is true on rows $(\TT,\FF), (\FF,\TT), (\FF,\FF)$, so
\[
p \to \neg q \;\equiv\;
(p \wedge \neg q) \vee (\neg p \wedge q) \vee (\neg p \wedge \neg q).
\]
\end{solution}

\begin{solution}{3.6}
$p \oplus q$ is false on rows $(\TT,\TT)$ and $(\FF,\FF)$. Each false row
yields the maxterm that negates it:
$(\neg p \vee \neg q)$ for $(\TT,\TT)$ and $(p \vee q)$ for $(\FF,\FF)$.
Hence $p \oplus q \equiv (p \vee q) \wedge (\neg p \vee \neg q)$.
\end{solution}

\begin{solution}{3.7}
Odd-parity rows: exactly one true ---
$(\TT,\FF,\FF), (\FF,\TT,\FF), (\FF,\FF,\TT)$ --- or all three true
$(\TT,\TT,\TT)$. Full DNF:
\[
(p\wedge\neg q\wedge\neg r) \vee (\neg p\wedge q\wedge\neg r) \vee
(\neg p\wedge\neg q\wedge r) \vee (p\wedge q\wedge r).
\]
Since $\oplus$ is associative and $x \oplus y$ is true when exactly one of
$x,y$ is true, $(p \oplus q) \oplus r$ is true when the number of true atoms
among $p,q,r$ is odd --- the same four rows. Hence
$f \equiv p \oplus q \oplus r$.
\end{solution}

\begin{solution}{3.8}
Distribute:
\[
(p \wedge q) \vee (r \wedge s)
\equiv (p \vee r) \wedge (p \vee s) \wedge (q \vee r) \wedge (q \vee s).
\]
With $k$ two-literal conjunctions $(a_1\wedge b_1)\vee\cdots\vee(a_k\wedge b_k)$,
distributing chooses one literal from each conjunct in every clause: $2^k$
clauses, and in general (for distinct variables) none is redundant.
\end{solution}

\begin{solution}{3.9}
Expand the CNF:
$(p \vee q) \wedge (p \vee r) \wedge (q \vee r)$. If at least two of $p,q,r$
are true, each clause contains a true literal, so the CNF is true. If at most
one is true, some pair among $\{p,q\},\{p,r\},\{q,r\}$ consists of two false
atoms, and the corresponding clause fails. So the CNF is true exactly when at
least two atoms are true --- the majority function $m$.
\end{solution}
